% Einheit 1 von 4 – Negative Zahlen kennen und ordnen (bank/rationale-zahlen/e1.jsonl, zusammenbau v0.1)
%% TODO Zweigzeile Teil 1: was hier gelernt wird (Satz in Schülersprache aus dem Einheitstitel) – Einheit 1
\einheitenkopf[e1]{Einheit 1 von 4 · Negative Zahlen kennen und ordnen}
\zweigzeile{ab Kl. 6, je nach Buch bis Kl. 7 · P10}
\setcounter{aufgabe}{7}

%% TODO Titel: Ich-kann-Satz fehlt in der Bank (Platzhalter: Kettenname) – Nr. 8
%% TODO Anweisung: ein Satz über der ersten Teilaufgabe fehlt in der Bank – Nr. 8
\begin{aufgabe}{Ordnen}
\swa{Schreibe an jeden Strich seine Zahl. Abstand zwischen zwei Strichen: $1$.}{\zahlenstrahl[xmin=-5,xmax=5,xstep=1]{0/0}}
\swa{Trage ein: $-4$, $3$ und $-8$.}{\zahlenstrahl[xmin=-10,xmax=10,xstep=1]{}}
\swa{Trage ein: $-5$, $1$ und $-9$.}{\zahlenstrahl[xmin=-10,xmax=10,xstep=1]{}}
\swa{Trage ein: $8$, $-2$ und $-9$.}{\zahlenstrahl[xmin=-10,xmax=10,xstep=1]{}}
\swa{Temperaturen: $-6$ °C, $2$ °C, $-1$ °C – trage sie ein.}{\zahlenstrahl[xmin=-10,xmax=10,xstep=1]{}}
\swa{Kontostände: $-7$ €, $4$ €, $-4$ € – trage sie ein.}{\zahlenstrahl[xmin=-10,xmax=10,xstep=1]{}}
\swfrage{Was bleibt gleich, was ändert sich?}
%% TODO Darstellung für schwach fehlt in der Bank (rationale-zahlen-e1-k1-s2-v1; Abschnitt „Für schwache Schüler“)
\swz{$-17$ – Gegenzahl und Betrag?}{}
%% TODO Darstellung für schwach fehlt in der Bank (rationale-zahlen-e1-k1-s3-v1; Abschnitt „Für schwache Schüler“)
\swz{$-15$ oder $-9$ – welche ist kleiner?}{}
\swa{Trage ein: $-2{,}5$, $0{,}5$ und $-\frac{3}{2}$.}{\zahlenstrahl[xmin=-3,xmax=3,xstep=0.5,karo=1]{}}
%% TODO Darstellung für schwach fehlt in der Bank (rationale-zahlen-e1-k1-s5-v1; Abschnitt „Für schwache Schüler“)
\swz{Ordne aufsteigend: $-0{,}8$; $\frac{1}{4}$; $-1{,}2$; $-\frac{3}{4}$}{}
\swz[3]{Gib eine Zahl an, die größer ist als $-320$. \hfill (P10 2015 OS)}{}
\swz[3]{Ordne aufsteigend: $-\frac{2}{5}$; $1{,}3$; $-0{,}43$; $0{,}9$ \hfill (P10 2014 OS)}{}
\end{aufgabe}

%% TODO Titel: Ich-kann-Satz fehlt in der Bank (Platzhalter: Kettenname) – Nr. 9
%% TODO Anweisung: ein Satz über der ersten Teilaufgabe fehlt in der Bank – Nr. 9
\begin{aufgabe}{Mitte zweier Zahlen}
%% TODO Darstellung für schwach fehlt in der Bank (rationale-zahlen-e1-k2-s1-v1; Abschnitt „Für schwache Schüler“)
\swz{Mitte von $-8$ und $2$?}{}
\end{aufgabe}

%% TODO Titel: Ich-kann-Satz fehlt in der Bank (Platzhalter: Kettenname) – Nr. 10
%% TODO Anweisung: ein Satz über der ersten Teilaufgabe fehlt in der Bank – Nr. 10
\begin{aufgabe}{runden}
%% TODO Darstellung für schwach fehlt in der Bank (rationale-zahlen-e1-k3-s1-v1; Abschnitt „Für schwache Schüler“)
\swz{$-3{,}47$ auf Zehntel runden}{}
\end{aufgabe}

%% TODO Titel: Ich-kann-Satz fehlt in der Bank (Platzhalter: Kettenname) – Nr. 11
%% TODO Anweisung: ein Satz über der ersten Teilaufgabe fehlt in der Bank – Nr. 11
\begin{aufgabe}{Punkte in vier Quadranten}
\swz{Punkt P – welche Koordinaten? Lies ab.}{\begin{ksys}[xmin=-5,xmax=5,ymin=-5,ymax=5,ablesen] \punkt{-3}{2}{P} \end{ksys}}
\end{aufgabe}

%% TODO Titel: Ich-kann-Satz fehlt in der Bank (Platzhalter: Kettenname) – Nr. 12
%% TODO Anweisung: ein Satz über der ersten Teilaufgabe fehlt in der Bank – Nr. 12
\begin{aufgabe}{Ordnen – Fehler finden, Begründen, Darstellungswechsel, Anwendung}
\swz[3]{Tim schreibt: $-45 > -30$, weil $45 > 30$. Finde den Fehler.}{}
\swfrage{Warum ist $-17$ kleiner als $-4$? Begründe.}
\swz{Ein Taucher ist $18$ m unter dem Meeresspiegel. Welche Zahl?}{}
\swz[3]{Tiefste Punkte an Land: Kaspisches Meer $-28$ m, Totes Meer $-430$ m, Death Valley $-86$ m, Qattara-Senke $-133$ m. Welcher Ort liegt am tiefsten?}{}
\end{aufgabe}

\uebersichtskasten{Zahlengerade: links wird es kleiner. Jede negative Zahl ist kleiner als jede positive; von zwei negativen ist die mit dem größeren Betrag die kleinere. \\ −7 < −2 < 0 < 1,5 \quad Gegenzahl von −2 ist 2; Betrag |−7| = 7}
